\chapter*{Заключение}						% Заголовок
\addcontentsline{toc}{chapter}{Заключение}	% Добавляем его в оглавление

%% Согласно ГОСТ Р 7.0.11-2011:
%% 5.3.3 В заключении диссертации излагают итоги выполненного исследования, рекомендации, перспективы дальнейшей разработки темы.
%% 9.2.3 В заключении автореферата диссертации излагают итоги данного исследования, рекомендации и перспективы дальнейшей разработки темы.
%% Поэтому имеет смысл сделать эту часть общей и загрузить из одного файла в автореферат и в диссертацию:

По итогам проделанной работы можно сделать следующие выводы:

\textbf{В главе 1} проведен анализ существующих решений в области разработки систем подвижности тренажеров авто- и спецтехники. Формализованы типовые эффекты движения ТС. Сформулированы требования к системе подвижности.

\textbf{В главе 2} получены уравнения обратной и прямой задач кинематики, рассмотрено решение ПЗК аналитическим и численным методами. Проведено исследование точности приближенного аналитического решения ПЗК. Была также получена зона безопасных положений платформы, вытекающая из геометрических размеров платформы. 

\textbf{В главе 3} была сформулирована модель динамики системы подвижности с тремя степенями свободы и кривошипно-шатунным 
механизмом привода звеньев. Получена рабочая зона, вытекающая из паспортных значений вращающих моментов мотор-редукторов
Параметры полученной модели были уточнены по результатам эксперимента на реальном оборудовании, что позволило рассчитывать вращающие моменты по заданным движениям платформы с приемлемой точностью.  

\textbf{В главе 4} приведены исследования алгоритмов управления СП. Рассмотрены два основных подхода к построению системы управления, произведено сравнение характеристик СП по критерию среднего квадрата ошибки.

Результаты работы могут быть обобщены на другие типы кинематических схем систем подвижности и использоваться в алгоритмах управления при необходимости учета динамических характеристик объекта, закрепленного на подвижном основании СП.
